\section{Prime estimates}\label{sec:primes}

The kernel construction requires two uniform estimates: an asymptotic
formula for primes in residue classes, and cancellation in polynomial
exponential sums over primes.  We derive the first from a zero-free
half-plane theorem and prove the second by Vaughan's identity and Weyl
differencing.  The latter argument also records why the product cutoff
and the arithmetic coefficients do not obstruct mixed differencing.

For positive integers $n$, put
\[
 \Lambda'(n)=(\log n)\mathbf 1_{\{n\text{ is prime}\}},\qquad
 \Lambda(n)=
 \begin{cases}\log p,&n=p^j\text{ for a prime }p\text{ and }j\ge1,\\
 0,&\text{otherwise}.
 \end{cases}
\]
All sums below run over integers, including sums whose endpoints are
real.  The constants in this section are independent of the lower
coefficients of the polynomial phases.

\subsection{Primes in arithmetic progressions}

Here is the precise analytic input from the companion zero-free theorem.
Only its assertion about Dirichlet $L$-functions is needed.

\begin{theorem}[Zero-free input, {\cite[Theorem~1.1]{ZeroFree}}]
\label{thm:zero-free-input}
For every integer $b\ge1$ and every Dirichlet character $\chi$ modulo
$b$, the Dirichlet $L$-function $L(s,\chi)$ has no zero in
$\Re s>7/8$.  This includes imprimitive characters and the trivial
character of modulus $1$, whose $L$-function is $\zeta(s)$.  A pole at
$s=1$ for a principal character is allowed.
\end{theorem}

\begin{proposition}[Uniform distribution in residue classes]
\label{prop:prime-progressions}
Uniformly for $Y\ge2$, $0\le x\le Y$, integers $b\ge1$, and reduced
residue classes $c$ modulo $b$,
\begin{equation}\label{eq:prime-progression}
 \sum_{\substack{n\le x\\n\equiv c\pmod b}}\Lambda'(n)
 =\frac{x}{\varphi(b)}
   +O\bigl(Y^{31/32}\log^2(2bY)\bigr).
\end{equation}
For a nonreduced class modulo $b>1$, the sum is $O(\log b)$ and has
zero main term.
\end{proposition}

\begin{proof}
Set $\sigma_0=15/16$.  We first prove the logarithmic derivative estimate
\begin{equation}\label{eq:log-derivative}
 \frac{L'}{L}(s,\chi)\ll\log\bigl(2b(2+|\Im s|)\bigr)
\end{equation}
on $\Re s=\sigma_0$, and throughout $\sigma_0\le\Re s\le2$ when
$|\Im s|\ge1$.  The constants are absolute, including for imprimitive
characters.

Write $\delta_\chi=1$ for a principal character and $0$ otherwise, and
let
\[
 G_\chi(s)=(s-1)^{\delta_\chi}L(s,\chi).
\]
Fix radii
$2-\sigma_0<r<R<2-7/8$ and use disks of radius $R$ centered at
$s_t=2+it$.  Theorem~\ref{thm:zero-free-input} makes $G_\chi$
holomorphic and nonvanishing on each disk, with the principal pole
removed.  To bound its size there, use periodicity to write
\[
 \sum_{n\le u}\chi(n)=\overline\chi\,u+E_\chi(u),\qquad
 |E_\chi(u)|\le2b,\qquad
 \overline\chi=\begin{cases}\varphi(b)/b,&\delta_\chi=1,\\0,&\delta_\chi=0.
 \end{cases}
\]
Partial summation continues the Dirichlet series to $\Re s>0$ as
\[
 L(s,\chi)=\overline\chi\,\frac{s}{s-1}
       +s\int_1^\infty E_\chi(u)u^{-s-1}\,du.
\]
Consequently, on the disk,
$|G_\chi(s)|\le C[2b(2+|t|)]^C$.  At its center the absolutely
convergent Euler product gives $|L(2+it,\chi)|\ge\zeta(4)/\zeta(2)$;
hence $|G_\chi(s_t)|$ is bounded below by the same positive constant.
Choose a holomorphic logarithm and subtract its value at $s_t$.
Its real part is bounded above by
$C\log(2b(2+|t|))$.  Borel--Carath\'eodory on an intermediate disk,
followed by Cauchy's estimate on the disk of radius $r$, now gives
\[
 \frac{G_\chi'}{G_\chi}(s)
 \ll\log(2b(2+|t|))\qquad (|s-s_t|\le r).
\]
Subtracting $\delta_\chi/(s-1)$ proves
\eqref{eq:log-derivative} in the asserted ranges.

We next pass through a smoothed sum so that the contour integral is
absolutely convergent.  Mellin inversion for $(1-u)_+$ gives, for $x\ge1$,
\begin{equation}\label{eq:smoothed-character-primes}
 \sum_{n\le x}(x-n)\Lambda(n)\chi(n)
 =\frac{1}{2\pi i}\int_{(2)}
   -\frac{L'}{L}(s,\chi)\frac{x^{s+1}}{s(s+1)}\,ds.
\end{equation}
Here the interchange with the Dirichlet series is absolute.  Move the
line to $\Re s=\sigma_0$.  Estimate~\eqref{eq:log-derivative} bounds
the horizontal integrals by a constant times
$x^3\log(2b(2+T))/T^2$, which tends to zero with the height $T$.
The sole crossed pole is at $s=1$ for the principal character; its
residue is $x^2/2$.  On the new vertical line, integration of
$\log(2b(2+|t|))/(1+t^2)$ gives
\begin{equation}\label{eq:smoothed-character-result}
 \sum_{n\le x}(x-n)\Lambda(n)\chi(n)
 =\delta_\chi\frac{x^2}{2}
      +O\bigl(x^{1+\sigma_0}\log(2b)\bigr).
\end{equation}

For $(c,b)=1$, character orthogonality yields
\[
 A_c(x):=\sum_{\substack{n\le x\\n\equiv c\pmod b}}
                 (x-n)\Lambda(n)
 =\frac{x^2}{2\varphi(b)}
      +O\bigl(x^{1+\sigma_0}\log(2b)\bigr).
\]
There are $\varphi(b)$ character errors and a factor $1/\varphi(b)$,
so the error is uniform in $b$.  Crucially, sharp differencing is now
performed on this nonnegative progression sum, rather than on a complex
character sum.  For $0<H<x$, positivity of $\Lambda$ gives
\[
 \frac{A_c(x)-A_c(x-H)}{H}
 \le \sum_{\substack{n\le x\\n\equiv c\pmod b}}\Lambda(n)
 \le \frac{A_c(x+H)-A_c(x)}{H}.
\]
Take $H=x^{31/32}$ for sufficiently large $x$.  The main terms differ
from $x/\varphi(b)$ by $O(H)$, and the errors are
\[
 O\bigl(x^{1+\sigma_0}/H\cdot\log(2b)\bigr)
   =O\bigl(x^{31/32}\log(2b)\bigr).
\]
Bounded $x$ is harmless.  Finally,
\[
 \sum_{\substack{p^j\le Y\\j\ge2}}\log p
 \ll Y^{1/2}\log^2(2Y),
\]
so removing higher prime powers proves~\eqref{eq:prime-progression}.
If a prime lies in a nonreduced residue class modulo $b$, it divides
$b$; the sum of the logarithms of distinct such primes is at most
$\log b$.  This proves the last assertion.
\end{proof}

\subsection{A minor-arc estimate over primes}

\begin{proposition}[Prime polynomial sums]\label{prop:prime-minor-arcs}
For fixed $k\ge2$ and $0<\rho<1$, there is a constant
$c_m=c_m(k,\rho)\in(0,1/4)$ with the following property.  If $Y$ is
sufficiently large, $I\subseteq[1,Y]$ is an interval, and
$f\in\R[x]$ has degree at most $k$ with coefficient $b_k$ of $x^k$,
then
\begin{equation}\label{eq:prime-minor-arc}
 \left|\sum_{n\in I}\Lambda'(n)e(f(n))\right|
 \ll_{k,\rho}Y^{1-c_m},
\end{equation}
provided $\theta=(k!)^2b_k$ has an approximation $a/q$ satisfying
\begin{equation}\label{eq:minor-rational-range}
 (a,q)=1,\qquad |\theta-a/q|\le q^{-2},\qquad
 Y^\rho\le q\le Y^{k-\rho}.
\end{equation}
\end{proposition}

Vaughan's identity will split the prime weight into sums with one short
variable and bilinear sums.  Weyl differencing turns each into averages
of linear phases, which rational approximation then bounds.  We first
prove the needed spacing and differencing estimates.  Uniformity in the
lower coefficients and in the interval will later permit additive linear
twists for periodic restrictions and partial summation for smooth weights
in the kernel construction.

Write $\|u\|$ for the distance of $u\in\R$ to the
nearest integer, and interpret $\min(M,\|u\|^{-1})$ as $M$ when
$\|u\|=0$.

\begin{lemma}[Rational spacing]\label{lem:rational-spacing}
Suppose $(a,q)=1$, $q\ge2$, and $|\theta-a/q|\le q^{-2}$.
For $R\ge1$ and $M\ge1$,
\begin{equation}\label{eq:rational-spacing}
 \sum_{1\le r\le R}\min(M,\|\theta r\|^{-1})
 \ll (R/q+1)\bigl(M+q\log(2q)\bigr).
\end{equation}
\end{lemma}

\begin{proof}
Within a block of $\lfloor q/2\rfloor$ consecutive integers, two distinct
indices differ by a nonzero $d$ with $|d|<q/2$.  Thus
\[
 \|\theta d\|\ge \|ad/q\|-|d|/q^2\ge1/(2q).
\]
The associated points are $1/(2q)$ separated on the circle.  The nearest
point to zero contributes at most $M$, and summing the others in bands
of width $1/(2q)$ gives $O(q\log(2q))$, with another $O(M)$ for the
bounded number of points in the first band.  Partitioning the indices
into such blocks proves the result, also when $q=2$ or $3$.
\end{proof}

We will repeatedly group a fixed number of nonzero integer factors by
their product.  If that product has absolute value at most $Y^{C_k}$,
the number of such factorizations is $Y^{o(1)}$, uniformly in the
product.  Indeed, its prime factorization bounds the number by a fixed
order divisor function, up to a bounded number of sign choices.  For
every $\epsilon>0$ that function is $O_{k,\epsilon}(Y^\epsilon)$:
allocating a prime exponent among finitely many factors gives a
polynomial in the exponent, which is bounded by any fixed positive
power of the prime power after allowing a constant for finitely many
small primes.

The following two differencing bounds isolate the treatment of
coefficients and summation domains.  They are forms of classical Weyl
differencing; see \cite{Weyl1916,Vaughan1997} for the method.

\begin{lemma}[One-variable and mixed differencing]
\label{lem:prime-differencing}
Let $k\ge2$, and let $f\in\R[x]$ have degree at most $k$, with
coefficient $b_k$ of $x^k$.

If $J$ is an interval of integers of length at most an integer $M\ge1$,
and $m$ is a positive integer, then
\begin{equation}\label{eq:type-one-differencing}
 \left|\frac1M\sum_{n\in J}e(f(mn))\right|^{2^{k-1}}
 \ll_k\frac1{M^k}
 \sum_{|u_1|,\ldots,|u_{k-1}|<M}
 \min\bigl(M,\|k!b_km^ku_1\cdots u_{k-1}\|^{-1}\bigr).
\end{equation}

Next let $A,B$ be positive integers, let $I\subseteq[1,Y]$ be an
interval, and let $|a_m|,|d_n|\le1$.  Set
\[
 S=\sum_{\substack{A\le m<2A\\B\le n<2B}}
       a_md_n\mathbf 1_{\{mn\in I\}}e(f(mn)),\qquad
 \theta=(k!)^2b_k.
\]
Then
\begin{equation}\label{eq:type-two-differencing}
 \left|\frac{S}{AB}\right|^{2^{2k-1}}
 \ll_k\frac1{A^kB^k}
 \sum_{\substack{|u_1|,\ldots,|u_k|<B\\
                  |v_1|,\ldots,|v_{k-1}|<A}}
 \min\bigl(A,\|\theta u_1\cdots u_kv_1\cdots v_{k-1}\|^{-1}\bigr).
\end{equation}
Zero shifts are included in both sums.
\end{lemma}

\begin{proof}
We give the mixed argument explicitly.  For a function $g$ of two
integer variables, define
\[
 \Delta^n_u g(m,n)=g(m,n+u)-g(m,n),\qquad
 \Delta^m_v g(m,n)=g(m+v,n)-g(m,n).
\]
Let $D_0$ be the set of pairs in the displayed box with $mn\in I$.
Extend $a_m$ and $d_n$ by zero outside their respective intervals,
and extend every summand by zero outside $D_0$.  Cauchy--Schwarz first
removes $a_m$:
\[
 \left|\frac{S}{AB}\right|^2
 \le\frac1A\sum_{A\le m<2A}
       \left|\frac1B\sum_n d_n\mathbf 1_{D_0}(m,n)e(f(mn))\right|^2.
\]
Expanding the square and grouping by $u=n'-n$ bounds its right side by
\[
 \frac1B\sum_{|u|<B}
 \left|\frac1{AB}\sum_{m,n}
    d_{n+u}\overline{d_n}\,
    \mathbf 1_{D_0}(m,n+u)\mathbf 1_{D_0}(m,n)
    e\bigl(\Delta^n_u f(mn)\bigr)\right|.
\]
For fixed $u$ the new weight depends only on $n$ and has modulus at
most one.  This is exactly the form needed for another $n$-difference.
In particular, the product cutoff remains inside the sum; it has not
been replaced by a rectangular domain.

For $j$ fixed shifts $\mathbf u=(u_1,\ldots,u_j)$, let
\[
 \begin{split}
 D_{\mathbf u}&=\{(m,n):
      (m,n+\varepsilon\cdot\mathbf u)\in D_0
                   \text{ for every }\varepsilon\in\{0,1\}^j\},\\
 F_{\mathbf u}(m,n)&=
    \Delta^n_{u_j}\cdots\Delta^n_{u_1}f(mn).
 \end{split}
\]
The accompanying coefficient $d_{\mathbf u}(n)$ is defined recursively
by $d_\varnothing(n)=d_n$ and
\[
 d_{(\mathbf u,u)}(n)
   =d_{\mathbf u}(n+u)\overline{d_{\mathbf u}(n)}.
\]
It depends only on $n$ and is bounded by one.  At each subsequent step,
Cauchy--Schwarz in the existing shifts costs only a constant depending
on $j$, because their total normalized mass is at most $2^j$.
Cauchy--Schwarz in $m$, followed by the same square expansion, adds one
$n$-shift.  Induction therefore gives
\begin{equation}\label{eq:n-difference-stage}
 \left|\frac{S}{AB}\right|^{2^k}
 \ll_k\frac1{B^k}\sum_{|u_1|,\ldots,|u_k|<B}
 \left|\frac1{AB}\sum_{m,n}
       d_{\mathbf u}(n)\mathbf 1_{D_{\mathbf u}}(m,n)
       e(F_{\mathbf u}(m,n))\right|.
\end{equation}
All summands still have $m$ and $n$ in their original intervals, as the
zero vertex of every translated cube belongs to $D_0$.

After $k$ differences in $n$, every lower-degree term vanishes, and
\[
 F_{\mathbf u}(m,n)=k!b_k u_1\cdots u_k m^k.
\]
For $0\le j\le k-1$ and $m$-shifts
$\mathbf v=(v_1,\ldots,v_j)$, define the domains
\[
 \begin{split}
 D_{\mathbf u,\mathbf v}=\{(m,n):\;&
 (m+\eta\cdot\mathbf v,n+\varepsilon\cdot\mathbf u)\in D_0\\
 &\text{for every }\eta\in\{0,1\}^j,
                         \ \varepsilon\in\{0,1\}^k\}.
 \end{split}
\]
Thus $D_{\mathbf u,\varnothing}=D_{\mathbf u}$.  For $1\le j\le k-1$, let
\[
 T_{\mathbf u,\mathbf v}
 =\frac1{AB}\sum_{m,n}\mathbf 1_{D_{\mathbf u,\mathbf v}}(m,n)
 e\bigl(\Delta^m_{v_j}\cdots\Delta^m_{v_1}F_{\mathbf u}(m,n)\bigr).
\]
To enter this unweighted family, remove $d_{\mathbf u}$ by
Cauchy--Schwarz in $n$:
\[
 \left|\frac1{AB}\sum_{m,n}
 d_{\mathbf u}(n)\mathbf 1_{D_{\mathbf u}}(m,n)e(F_{\mathbf u}(m,n))
 \right|^2
 \le\frac1B\sum_{B\le n<2B}
 \left|\frac1A\sum_m
    \mathbf 1_{D_{\mathbf u}}(m,n)e(F_{\mathbf u}(m,n))\right|^2.
\]
Expanding the square bounds this by
$A^{-1}\sum_{|v_1|<A}|T_{\mathbf u,(v_1)}|$.
This is the first $m$-differencing step, so removing the coefficients
requires no additional squaring.  For $1\le j\le k-2$, Cauchy--Schwarz
in $n$ and the same square expansion give
\[
 |T_{\mathbf u,\mathbf v}|^2
 \le\frac1A\sum_{|v|<A}|T_{\mathbf u,(\mathbf v,v)}|.
\]
Applying Cauchy--Schwarz also to all existing shift variables, whose
total normalized mass is bounded in terms of $k$, now yields
\begin{equation}\label{eq:m-difference-stage}
 \left|\frac{S}{AB}\right|^{2^{k+j}}
 \ll_k\frac1{B^kA^j}
 \sum_{\substack{|u_1|,\ldots,|u_k|<B\\|v_1|,\ldots,|v_j|<A}}
          |T_{\mathbf u,\mathbf v}|\qquad(1\le j\le k-1).
\end{equation}
This accounts for all $k+(k-1)$ squarings.

Take $j=k-1$.  For fixed $n$, the section of $D_{\mathbf u,\mathbf v}$
is an interval of $m$'s, possibly empty.  In fact each
condition puts $m+\eta\cdot\mathbf v$ in $[A,2A)$ and its product with
the positive number $n+\varepsilon\cdot\mathbf u$ in $I$; it therefore
specifies an interval in $m$.  Their intersection is an interval of
length at most $A$.  The final phase
$\Delta^m_{v_{k-1}}\cdots\Delta^m_{v_1}F_{\mathbf u}$ is linear in $m$
with coefficient
\[
 (k!)^2b_k\,u_1\cdots u_kv_1\cdots v_{k-1}.
\]
A geometric-sum estimate bounds each $m$-sum by the corresponding
minimum in~\eqref{eq:type-two-differencing}.  Summing over at most $B$
values of $n$, the remaining normalized sum is bounded by that minimum
divided by $A$.  The shifts have normalizing factor
$B^{-k}A^{-(k-1)}$, which proves~\eqref{eq:type-two-differencing}.

For~\eqref{eq:type-one-differencing}, perform the same square expansion
$k-1$ times in a single variable.  The translated intersections of $J$
are intervals of length at most $M$, and the final linear coefficient
is $k!b_km^ku_1\cdots u_{k-1}$.  The normalizing factor is
$M^{-(k-1)}$, followed by $M^{-1}$ for the final geometric sum.
This proves the first bound as well.
\end{proof}

\begin{proof}[Proof of Proposition~\ref{prop:prime-minor-arcs}]
We now apply the differencing bounds to Vaughan's decomposition.
We first prove the estimate with $\Lambda$ in place of $\Lambda'$.
Put $e_0=\rho/(100k)$ and $Z=Y^{e_0}$.  For an arithmetic function
$g$, write $g_{\le Z}(n)=g(n)\mathbf 1_{\{n\le Z\}}$ and
$g_{>Z}=g-g_{\le Z}$.  Let $\mu$ denote the M\"obius function, let
$1$ denote the constant arithmetic function, and use $*$ for Dirichlet
convolution.  Vaughan's identity, in the form used here, is
\begin{equation}\label{eq:vaughan-identity}
 \Lambda=\Lambda_{\le Z}+\mu_{\le Z}*\log
       -\mu_{\le Z}*\Lambda_{\le Z}*1
       +\mu_{>Z}*1*\Lambda_{>Z}.
\end{equation}
This decomposition is classical; see \cite[Chapter~3]{Vaughan1997}.
It follows directly from $\log=1*\Lambda$ and
$\mu*1=\mathbf 1_{\{1\}}$: split $\mu$ at $Z$ in $\mu*\log$, and then
split $\Lambda$ at $Z$ in the term $\mu_{>Z}*1*\Lambda$.

\emph{The two sums with a short variable.}
The middle two terms of~\eqref{eq:vaughan-identity} reduce to sums
with an outer variable $m\le Z^2$ and an inner variable $n\le Y/m$.
The outer coefficient is either $\mu(m)$ or
$\sum_{ab=m,\,a,b\le Z}\mu(a)\Lambda(b)$, and therefore has size
$Y^{o(1)}$.  The inner weight is respectively $\log n$ or $1$.
Fix $m$, let $M=\lceil Y/m\rceil$, and consider any interval subsum
$S_m=\sum_{n\in J}e(f(mn))$ of the inner sum.

Apply~\eqref{eq:type-one-differencing}.  For every integer $r$,
\[
 \min(M,\|k!b_kr\|^{-1})
 \le k!\min(M,\|\theta r\|^{-1}),
\]
since $\|k!u\|\le k!\|u\|$.  Tuples with a zero shift contribute
$O_k(M^{-1})$ after normalization.  For the others group by
$r=m^k|u_1\cdots u_{k-1}|\le m^kM^{k-1}$ and apply
Lemma~\ref{lem:rational-spacing}.  The divisor bound gives
\begin{equation}\label{eq:type-one-saving}
 \left|\frac{S_m}{M}\right|^{2^{k-1}}
 \ll Y^{o(1)}\left(
 \frac1M+\frac{m^k}{q}+\frac{m^k}{M}
       +\frac1{M^{k-1}}+\frac q{M^k}\right).
\end{equation}
Here and below logarithmic factors are included in $Y^{o(1)}$.
For completeness, the four terms after $1/M$ result from expanding
\[
 M^{-k}\left(\frac{m^kM^{k-1}}q+1\right)(M+q\log(2q)).
\]
As $m\le Y^{2e_0}$ and $M\ge Y^{1-2e_0}$,
\[
 \frac{m^k}{q},\ \frac q{M^k}\le Y^{-\rho+2ke_0},\qquad
 \frac{m^k}{M}\le Y^{-1+2(k+1)e_0}.
\]
The remaining terms also save a positive power of $Y$.  Taking the
$2^{k-1}$-st root in~\eqref{eq:type-one-saving} therefore gives
$S_m\ll M Y^{-c}$ for some $c=c(k,\rho)>0$, uniformly in the subinterval.
Partial summation permits the weight $\log n$ with a further
$O(\log Y)$ factor.  Since
\[
 \sum_{m\le Z^2}\lceil Y/m\rceil\ll Y\log(2Y),
\]
both short-variable terms save a fixed power of $Y$ after including
the outer coefficients.

\emph{The bilinear sum.}
For the final term of~\eqref{eq:vaughan-identity}, group the factors as
$(\mu_{>Z}*1)(m)\Lambda_{>Z}(n)$.  Each coefficient is $Y^{o(1)}$
when $mn\le Y$, and a nonzero coefficient requires both $m>Z$ and
$n>Z$.  Partition into dyadic boxes
\[
 A\le m<2A,\qquad B\le n<2B,
 \qquad A,B\ge Z/2,\qquad AB\le Y.
\]
There are $O(\log^2(2Y))$ relevant boxes.  Divide out the coefficient
bounds to reduce to $|a_m|,|d_n|\le1$.  If $AB<Y^{1-e_0}$,
the trivial estimate is already sufficient.  Otherwise put
$R=A^{k-1}B^k$ and apply~\eqref{eq:type-two-differencing}.

There are $k$ shifts of range $B$ and $k-1$ of range $A$.
Tuples containing a zero shift contribute
$O_k(A^{-1}+B^{-1})$.  Group the others by their nonzero absolute
product, at most $R$.  Lemma~\ref{lem:rational-spacing} and the divisor
bound imply
\begin{equation}\label{eq:type-two-saving}
 \left|\frac{S}{AB}\right|^{2^{2k-1}}
 \ll Y^{o(1)}\left(
 \frac1A+\frac1B+\frac1q+\frac1R+\frac q{AR}\right).
\end{equation}
Indeed the nonzero-shift bound before expansion is
\[
 \frac{Y^{o(1)}}{AR}(R/q+1)(A+q\log(2q));
\]
the additional term $O(\log(2q)/A)$ is included in $Y^{o(1)}/A$.
Now $A,B\ge Y^{e_0}/2$, while
\[
 AR=(AB)^k\ge Y^{k(1-e_0)},\qquad
 \frac q{AR}\le Y^{-\rho+ke_0}.
\]
Thus~\eqref{eq:type-two-saving} saves a fixed power.  Restoring the
coefficients and summing the boxes preserves a positive saving.

The remaining term $\Lambda_{\le Z}$ costs only $O(Z\log(2Z))$.
This proves the required estimate for $\Lambda$, with some positive
exponent depending only on $k,\rho$.  Removing higher prime powers
costs $O(Y^{1/2}\log^2(2Y))$, as in
Proposition~\ref{prop:prime-progressions}.  Decreasing the exponent
if necessary gives $c_m\in(0,1/4)$ and proves
\eqref{eq:prime-minor-arc}.
\end{proof}
